\documentclass[11pt,a4paper]{article}
\usepackage[utf8]{inputenc}
\usepackage[spanish]{babel}
\usepackage{amsmath, amssymb}
\usepackage[margin=2.5cm]{geometry}
\usepackage{fancyhdr}
\usepackage{tabularx}
\usepackage{booktabs}
\usepackage{tcolorbox}
\usepackage{enumitem}

% Configuración de colores corporativos UCB
\definecolor{UCBblue}{RGB}{0, 51, 153}
\definecolor{UCBgray}{RGB}{100, 100, 100}
\definecolor{UCBred}{RGB}{153, 0, 0}

% Configuración de encabezado y pie de página
\pagestyle{fancy}
\fancyhf{}
\fancyhead[L]{\textcolor{UCBblue}{\textbf{Ingeniería Mecatrónica - UCB Tarija}}}
\fancyhead[R]{\textcolor{UCBgray}{Guión de Clase - Sesión 3}}
\fancyfoot[C]{\thepage}
\renewcommand{\headrulewidth}{0.4pt}
\renewcommand{\headrule}{\hbox to\headwidth{\color{UCBblue}\leaders\hrule height \headrulewidth\hfill}}

\begin{document}

\begin{center}
    {\Large \textbf{\textcolor{UCBblue}{PLANIFICACIÓN DIDÁCTICA Y GUIÓN DE CLASE}}}\\
    \vspace{0.3cm}
    {\large \textbf{Sesión 3: Dinámica Frecuencial y Visualización Logarítmica}}\\
    \vspace{0.5cm}
    \textbf{Asignatura:} Taller de Nivelación - Álgebra Lineal Aplicada a la Teoría de Redes \\
    \textbf{Docente:} M.Sc. Ing. Bernardo Quiroga Turdera \quad \textbf{Duración:} 120 Minutos
\end{center}

\vspace{0.3cm}

\section*{1. Metadatos y Alineamiento Constructivo}
\begin{itemize}
    \item \textbf{Módulo:} 4 (Algoritmo Híbrido y Barrido Frecuencial).
    \item \textbf{Resultados de Aprendizaje (ABET 1 y 6):} Modelar dinámicas frecuenciales y generar visualizaciones de datos (Bode preliminar) automatizando el cálculo matricial en Python.
    \item \textbf{Estrategia Didáctica:} \textit{Just-In-Time Physics} (Reactancia vs. Frecuencia) y Aprendizaje Basado en Problemas (El ``límite humano'' del cálculo analítico).
    \item \textbf{Entregables Activos:} Diapositivas Beamer (S3), Guía Analítica del Estudiante (S3), Jupyter Notebook 04.
\end{itemize}

\section*{2. Cronograma de Ejecución (120 Minutos)}
\begin{table}[h!]
    \centering
    \renewcommand{\arraystretch}{1.3}
    \begin{tabularx}{\textwidth}{@{} c c X l @{}}
        \toprule
        \textbf{Hora} & \textbf{Fase CDIO} & \textbf{Actividad Principal} & \textbf{Soporte} \\
        \midrule
        17:00 - 17:15 & \textbf{Evaluar} & \textit{Warm-up}: Repaso de impedancias ($\mathbf{Z}$) y el uso del operador \texttt{1j}. & Pizarra \\
        17:15 - 17:40 & \textbf{Concebir} & La reactancia dinámica ($\omega$) y el porqué de la escala logarítmica (Décadas). & Beamer S3 \\
        17:40 - 18:00 & \textbf{Diseñar} & Misión ``El Límite Humano'': Cálculo a pulso de $2$ frecuencias extremas. & Guía (PDF) \\
        18:00 - 18:40 & \textbf{Implementar} & Funciones modulares (\texttt{def}) y vectores logarítmicos (\texttt{np.logspace}). & Notebook 04 \\
        18:40 - 19:00 & \textbf{Operar} & Dashboarding (\texttt{plt.semilogx}), detección de resonancia y cierre. & Notebook 04 \\
        \bottomrule
    \end{tabularx}
\end{table}

\section*{3. Desarrollo Detallado por Fases}

\subsection*{Fase 1: Fundamentación Teórica (17:15 - 17:40)}
\textit{Objetivo: Destruir el paradigma estático de la CC e introducir el dominio frecuencial.}
\begin{itemize}
    \item \textbf{Reactancia Dinámica:} Explicar físicamente por qué una bobina bloquea las altas frecuencias ($X_L = \omega L$) y el capacitor las deja pasar ($X_C = -1/\omega C$).
    \item \textbf{El Problema del Espectro:} Plantear el desafío industrial de sintonizar un filtro PID o evitar que una articulación robótica entre en resonancia destructiva. Esto requiere evaluar \textit{miles} de frecuencias, no solo una.
    \item \textbf{La Escala Logarítmica:} Justificar matemáticamente el uso de ejes semilogarítmicos para no comprimir las bajas frecuencias. Introducir el concepto de "Década".
\end{itemize}

\subsection*{Fase 2: Misión de Ingeniería - El Límite Humano (17:40 - 18:00)}
\textit{Objetivo: Generar una frustración intencional (constructiva) en el cálculo manual para valorar la programación.}
\begin{itemize}
    \item \textbf{Dinámica:} Entregar la Guía Analítica S3. Solicitar el cálculo manual de la magnitud de la corriente para $\omega = 100$ rad/s y $\omega = 10,000$ rad/s.
    \item \textbf{Intervención del Docente:} Cuando los estudiantes tarden más de 10 minutos en obtener apenas dos puntos del gráfico, detener la actividad. Explicar: \textit{"Acaban de experimentar el límite del cálculo analítico humano. Es por esto que los ingenieros programan"}.
\end{itemize}

\subsection*{Fase 3: Implementación Computacional y Dashboarding (18:00 - 19:00)}
\textit{Objetivo: Automatizar la física matricial mediante bucles y funciones.}
\begin{itemize}
    \item \textbf{Modularidad:} Guiar la programación de la función \texttt{def resolver\_red(w):}. Enseñar que la matriz $\mathbf{Z}$ ahora es dinámica y se reconstruye internamente para cada valor de $\omega$.
    \item \textbf{Vectores Logarítmicos:} Utilizar \texttt{np.logspace(1, 4, 500)} para generar 500 puntos de prueba y procesarlos con un bucle \texttt{for}.
    \item \textbf{Operación (Visualización):} Utilizar \texttt{plt.semilogx} para el despliegue del gráfico (Bode). Extraer computacionalmente la frecuencia donde la corriente es máxima (\texttt{np.argmax}).
\end{itemize}

\subsection*{Fase 4: Evaluación y Cierre Socrático}
\begin{tcolorbox}[colback=white,colframe=UCBgray,title=\textbf{Preguntas Socráticas de Cierre}]
\begin{enumerate}
    \item \textit{``En el gráfico logarítmico obtenido, observamos un 'pico' agudo (Resonancia). Físicamente, ¿qué le está ocurriendo a la reactancia del capacitor y la bobina en ese exacto punto de frecuencia para que la corriente se dispare?''} (Respuesta esperada: Se anulan mutuamente, $X_L = X_C$, dejando solo la resistencia pura).
    \item \textit{``¿Por qué un ingeniero mecatrónico necesitaría conocer esta frecuencia de resonancia antes de encender un brazo robótico?''}
\end{enumerate}
\end{tcolorbox}

\end{document}
